\tikzsetnextfilename{fig_ball_bouncing_perturbation_cycles_example}
%\resizebox{.5\textwidth}{!}{%


\begin{tikzpicture}[arrow/.style={thick, ->, >=latex}]
	%
	\centering
	\begin{axis}[%
		width=8cm,
		height=4.5cm,
		scale only axis,    
%				xmin=111.9,
%				xmax=114.3,
		xlabel={Temps},
%				ymin=-0.2,
		ymax=2.2,
		ylabel={Position verticale},
		axis lines=left,
%				xtick=\empty,
%				ytick=\empty,
		%axis line style={-latex},
		trim axis left, trim axis right,
		]
		%
		\addplot [
			color=Orient,
			solid,
			thick,
			]
			table [col sep=comma, x=t, y=xp] {misc/example_ball_bouncing_000.csv};
		%
		\addplot [
			color=Shiraz,
			solid,
			thick,
			]
			table [col sep=comma, x=t, y=xb] {misc/example_ball_bouncing_000.csv};
		%
		%
		\addplot [
		color=black,
		mark=asterisk,
		mark options={scale=1.75, line width=1pt},%
		]
		table[row sep=crcr]{
			142.927	0.2235\\
		};
		\coordinate (pert) at (axis cs: 142.927, 0.2235);
		%\tikzmark{end}
		%\node[cross,rotate=30,centered] at (axis cs:142.927, 0.2235) {};
		\draw[dashed, thick] (axis cs:\pgfkeysvalueof{/pgfplots/xmin}, 1.75) -- (axis cs:\pgfkeysvalueof{/pgfplots/xmax}, 1.75) node[at start, above, xshift=.75cm] {$h_2^*$};
		\node[above] at (axis cs:143.192, 1.85) {$c^{-1}$}; %1.6484
		\node[above] at (axis cs:144.326, 1.85) {$c^{1}$};  %1.8303
		\node[above] at (axis cs:145.455, 1.85) {$c^{2}$};  %1.7911
		\node[above] at (axis cs:146.652, 1.85) {$c^{3}$};  %1.7255
		\draw[dashed] (axis cs:142.647, \pgfkeysvalueof{/pgfplots/ymin}) -- (axis cs:142.647, \pgfkeysvalueof{/pgfplots/ymax});
		\draw[dashed] (axis cs:143.729, \pgfkeysvalueof{/pgfplots/ymin}) -- (axis cs:143.729, \pgfkeysvalueof{/pgfplots/ymax});
		\draw[dashed] (axis cs:144.892, \pgfkeysvalueof{/pgfplots/ymin}) -- (axis cs:144.892, \pgfkeysvalueof{/pgfplots/ymax});
		\draw[dashed] (axis cs:146.058, \pgfkeysvalueof{/pgfplots/ymin}) -- (axis cs:146.058, \pgfkeysvalueof{/pgfplots/ymax});
		\draw[dashed] (axis cs:147.211, \pgfkeysvalueof{/pgfplots/ymin}) -- (axis cs:147.211, \pgfkeysvalueof{/pgfplots/ymax});		
	\end{axis}
		\draw<2>[arrow] (pert) to[out=-70,in=110] ($(pert) + (-.25cm, -1.5cm)$) node [yshift=-.2cm, align=center] {\textbf<2>{\textcolor<2>{Prune}{perturbation}}};
\end{tikzpicture}
%}
%%
%